\section{Introduction}
\label{sec:intro}

The theory-based view of strategy holds that economic actors are theorists:
managers and entrepreneurs carry causal accounts of how value is created, and
those accounts---not merely their information or their incentives---direct
what they notice, which experiments they run, and what they do
\citep{FelinZenger2017,FelinGambardellaZenger2024}.  A decade of work has
established the consequences of holding a theory.  Theories guide
experimentation and search \citep{CamuffoEtAl2020,ChavdaGansStern2024},
support the acquisition of resources and rents \citep{EhrigZenger2024}, and
discipline entrepreneurial decision-making in the field
\citep{CamuffoEtAl2024}.  What this literature has not yet formalized is the
step before all of that: where a genuinely new theory comes from.  Existing
formal treatments begin with the theory present---as a menu of candidate
theories to select among \citep{CamuffoEtAl2024}, as a given causal
conjecture that directs search \citep{ChavdaGansStern2024}, or as a given
awareness advantage to be exploited competitively \citep{EhrigZenger2024}.
The origin of the theory is described verbally and documented empirically
\citep{OttHannah2024}, and the absence of formal machinery for it has been
explicitly noted as the field's open analytical problem
\citep{FelinGambardellaZenger2024}.

The difficulty is not a missing application of Bayesian reasoning; it is a
boundary of Bayesian reasoning.  A Bayesian agent updates beliefs over a
space of represented possibilities.  Conditioning reweights that space and
cannot enlarge it, so no sequence of updates makes an unrepresented causal
distinction available for belief or action.  Worse, when every represented
theory is false, the posterior does not sound an alarm: by the logic of
misspecified Bayesian learning \citep{Berk1966}, it concentrates on the least
wrong of the available theories, and the agent grows more confident, not
less.  Within-language updating is self-sealing.  The economics of growing
awareness has developed axiomatic accounts of how beliefs should extend when
the state space expands \citep{KarniViero2013,KarniViero2017,Viero2021}, but
in those models the expansion itself arrives exogenously; models of discovery
in games describe how play can reveal objectively available actions
\citep{Schipper2021} without explaining how experience generates a
substantive causal question or its answer.  What is missing, in both the
strategy literature and the economics of unawareness, is a model of the
transition: from a broken causal language, through a directed question, to a
new representation, and back to ordinary Bayesian life.

This paper models that transition, embedded in strategic interaction.  Its
organizing object is the chain
\[
\text{dissonance}
\longrightarrow
\text{inquiry}
\longrightarrow
\text{refined awareness}
\longrightarrow
\text{new models}
\longrightarrow
\text{Bayesian judgment},
\]
the process by which an agent whose causal representation has failed regains
the capacity to be a Bayesian.  The back end of the model follows the
subjective-rationality tradition \citep{KalaiLehrer1993b,Ryall2003}: a finite
set of agents adopt causal theories, combine them with beliefs and
conjectures to choose committed policies, and the policy profile determines
every agent's consequences through an objective causal system that absorbs
all market and environmental detail.  The front end is new.  Agents represent
the environment through awareness partitions of a state space; the theories
they can formulate assign elementary causal mechanisms to awareness cells;
and the partition itself---the language---is endogenous.  An awareness cell
is an implicit claim that the states it pools are causally alike.  Strategic
experience can refute that claim: when the statistical record rejects every
expressible theory, the agent faces \emph{causal dissonance}, a determinate
defect with no expressible repair.  Dissonance makes a question salient
without providing its answer.  A costly, fallible, truth-blind inquiry
technology may then produce a refinement of the partition---the formal
footprint of an insight---after which a bridge prior restores a probability
space and Bayesian judgment resumes on the expanded language.

Five sets of results give the chain its content.  First, boundary results
locate exactly where Bayesian competence stops: conditioning cannot expand
the language, and the misspecified posterior is self-sealing, so detecting
representational failure requires a faculty formally distinct from
updating.  Second, dissonance directs inquiry without determining it.  The
record restricts admissible repairs---every restorative refinement must
separate the evidence clusters that no single mechanism can accommodate---but
generally many minimal repairs remain, observationally equivalent worlds
demand different answers, and no truth-blind technology can be uniformly
correct.  Fallibility is the price of a representation that economizes on
distinctions.  Third, the value of a question is computable before its
answers are conceivable.  A leverage statistic built from retained
experience---invariant across every representation that could answer the
question---prices inquiry at exactly the moment the posterior is undefined,
and it tracks whether an answer could change action, not how much uncertainty
remains.  Questions die rationally: sustained success extinguishes them,
and defeat concentrates both the evidence and the incentive to inquire in
the losing firm.  Fourth, the return to Bayes is disciplined by a warrant
account: because any successful repair fits the experience that generated it
by construction, retrospective fit carries no information, and---when
inquiry is truth-blind---neither does the arrival of the insight itself.
Post-insight confidence therefore exceeds warrant by a computable gap, and an
agent holding a minimal repair generally cannot even formulate the comparison
that would measure its own gap.  Fifth, because theories cause actions and
actions generate the evidence, the population dynamics close the loop.  Play
converges to configurations in which every agent permanently inhabits a
Bayesian world---models fit, questions are dead, inquiry is priced below
cost.  We call such a configuration a \emph{self-confirming unawareness
equilibrium}, and its interest is epistemic rather than strategic: it can
preserve false theories, heterogeneous awareness, and permanent performance
differences, and it does so precisely when neither of the two corrective
forces---strategic variety inside an agent's categories, or an environment
harsh enough to punish misservice observably---is present.

The persistence results deserve emphasis because of what carries them.
Performance differences persist in this model not because an equilibrium
concept freezes them but because the process that would erase them has died
economically: the disadvantaged configuration generates no dissonance, or
generates no leverage, or prices its question below the cost of asking.
Strategic variety is therefore traced to its source.  Agents with identical
preferences, identical data-processing competence, and identical objective
opportunities end up with different causal languages, different theories, and
different payoffs because their experiences posed different questions, their
technologies produced different repairs, and their subsequent positions
extinguished further inquiry at different points.  This is a microfoundation
for heterogeneity that begins before probability does.

A deliberately small example runs through the paper---two firms, four
customer states, two product architectures---in which every object can be
computed exactly: the four minimal repairs of a broken language, the two that
survive a simplicity filter, the impossibility of choosing between them from
the data, the exact price of the question, the exact warrant gap, and the
knife-edge condition under which a market's shared false theory protects
itself forever.  The example is a laboratory, not the theory; each of its
results instantiates a general statement.

The paper proceeds as follows.  \Cref{sec:literature} locates the
contribution.  \Cref{sec:model} presents the model.
\Cref{sec:bayesian} develops the Bayesian layer and its boundary;
\cref{sec:dissonance} defines dissonance and exposure;
\cref{sec:inquiry} analyzes restoration, direction, and fallibility;
\cref{sec:economics} prices inquiry; \cref{sec:pivot} develops formulation,
warrant, and the return to Bayesian competence; and \cref{sec:dynamics}
closes the loop with the population dynamics.  \Cref{sec:discussion}
discusses implications and the research program.
